\section{Convex geometry and the minimum robust dimension}
\label{sec:threshold-geometry}

Theorem~\ref{thm:C10} reduces the all-diagonal problem to
\[
T_\alpha(\beta)
=
\min_{x\in\Delta_2^\circ}
F_{\alpha,\beta}(x),
\qquad
F_{\alpha,\beta}(x)
=
\frac{h_\alpha(B_\beta(x))}{x_1x_2x_3}.
\]
We now show that this variational problem is globally well behaved.

\subsection{Strict convexity and uniqueness}

Parameterize the simplex by
\[
x_1=\frac{e^{y_1}}{e^{y_1}+e^{y_2}+1},
\qquad
x_2=\frac{e^{y_2}}{e^{y_1}+e^{y_2}+1},
\qquad
x_3=\frac1{e^{y_1}+e^{y_2}+1}.
\]
Set
\[
L(y)=\log(e^{y_1}+e^{y_2}+1),
\]
\[
Q_\beta(y)
=
\log\!\left(
\beta_{12}e^{y_1+y_2}
+\beta_{13}e^{y_1}
+\beta_{23}e^{y_2}
\right),
\]
and
\[
\phi_\alpha(s)=\log h_\alpha(e^s).
\]
Then
\[
G_{\alpha,\beta}(y)
:=
\log F_{\alpha,\beta}(x(y))
=
\phi_\alpha(Q_\beta-2L)+3L-y_1-y_2.
\]

\begin{lemma}[Elasticity bounds]
\label{lem:elasticity-sub}
For \(2/3<\alpha<1\) and \(b>0\),
\[
0<
E_\alpha(b)
:=
\frac{b h_\alpha'(b)}{h_\alpha(b)}
<
\frac32,
\qquad
E_\alpha'(b)>0.
\]
Equivalently,
\[
0<\phi_\alpha'(s)<3/2,
\qquad
\phi_\alpha''(s)>0.
\]
\end{lemma}

\begin{proof}
With
\[
u=\cos(\alpha\pi/2),
\qquad
K=1-4u^2,
\]
parameterize the cubic boundary by
\[
b=r(Kr-2u),
\qquad
h=r^2(1+2ur),
\qquad
r>2u/K.
\]
Then
\[
E
=
\frac{b}{h}\frac{dh/dr}{db/dr}
=
\frac{(Kr-2u)(1+3ur)}
{(Kr-u)(1+2ur)}.
\]
The two factors
\[
\frac{Kr-2u}{Kr-u},
\qquad
\frac{1+3ur}{1+2ur}
\]
increase strictly with \(r\), while \(E\to0\) as \(r\downarrow2u/K\) and \(E\to3/2\) as \(r\to\infty\). Since \(b(r)\) is strictly increasing, the result follows.
\end{proof}

\begin{theorem}[Global strict logit convexity]
\label{thm:strict-convexity}
For every \(\beta>0\) and \(2/3<\alpha<1\), \(G_{\alpha,\beta}\) is globally strictly convex and coercive. Hence \(F_{\alpha,\beta}\) has a unique nondegenerate minimizer
\[
x^*(\alpha,\beta)\in\Delta_2^\circ,
\]
and \(T_\alpha(\beta)\) and \(x^*(\alpha,\beta)\) depend smoothly on \((\alpha,\beta)\).
\end{theorem}

\begin{proof}
Writing \(s=Q_\beta-2L\),
\[
\nabla^2G
=
\phi_\alpha''(s)\nabla s\nabla s^\top
+
\phi_\alpha'(s)\nabla^2Q_\beta
+
\bigl(3-2\phi_\alpha'(s)\bigr)\nabla^2L.
\]
Here \(\nabla^2Q_\beta\succeq0\), while \(\nabla^2L\succ0\) because \(L\) is the log-sum-exp of the three affinely independent vectors \((1,0),(0,1),(0,0)\) with positive softmax weights. Lemma~\ref{lem:elasticity-sub} makes every scalar coefficient nonnegative and the coefficient of \(\nabla^2L\) strictly positive. Thus
\[
\nabla^2G(y)\succ0
\qquad
\forall y\in\mathbb R^2.
\]
As \(\|y\|\to\infty\), \(x(y)\) approaches the simplex boundary, \(x_1x_2x_3\to0\), and the numerator remains strictly positive; hence \(G\to\infty\). Coercivity gives a unique global minimizer, positive definiteness gives nondegeneracy, and the implicit-function theorem yields smooth dependence.
\end{proof}

\subsection{Symmetric slice and canonical family}

\begin{corollary}[Fully symmetric slice]
\label{cor:symmetric-slice}
If
\[
\beta_{12}=\beta_{13}=\beta_{23}=b>0,
\]
then
\[
x^*=(1/3,1/3,1/3),
\qquad
T_\alpha(b,b,b)=27h_\alpha(b/3).
\]
For \(b=1\),
\[
T_\alpha(1,1,1)=1+R_3(\alpha)^3,
\qquad
R_3(\alpha)
=
\frac{\sin(\alpha\pi/2)}
{\sin(\alpha\pi/2-\pi/3)}.
\]
\end{corollary}

\begin{proof}
Permutation symmetry and uniqueness force \(x^*\) to be the simplex center. Substitution gives the first identity; the second simplifies to the \(n=3\) single-cycle threshold of Siami \cite{Siami2021}.
\end{proof}

Consider
\[
A_\gamma
=
\begin{pmatrix}
-1&0&-\gamma\\
\gamma&-1&0\\
0&\gamma&-1
\end{pmatrix},
\qquad
\gamma>0.
\]
It has
\[
\beta=(1,1,1),
\qquad
\kappa=1+\gamma^3.
\]

\begin{corollary}[Exact cyclic threshold]
\label{cor:cyclic-family}
For \(2/3<\alpha<1\),
\[
A_\gamma\in\mathcal F_\alpha^{(3)}
\iff
\gamma<R_3(\alpha),
\]
whereas
\[
A_\gamma\in\mathcal D_H^{(3)}
\iff
\gamma<2.
\]
Thus
\[
2<\gamma<R_3(\alpha)
\]
is an exact one-parameter genuinely fractional D-stable band.
\end{corollary}

\begin{proof}
The fractional statement follows from Theorem~\ref{thm:C10} and Corollary~\ref{cor:symmetric-slice}; the classical endpoint is
\[
1+\gamma^3<T_1(1,1,1)=9.
\]
\end{proof}

\subsection{Order dependence and the classical limit}

\begin{theorem}[Strict monotonicity]
\label{thm:alpha-monotonicity}
For fixed \(\beta>0\),
\[
\frac23<\alpha_1<\alpha_2\le1
\quad\Longrightarrow\quad
T_{\alpha_1}(\beta)>T_{\alpha_2}(\beta).
\]
\end{theorem}

\begin{proof}
For fixed \(a,b>0\), the cubic boundary value \(H_\theta(a,b)\) strictly decreases as the forbidden angle \(\theta\) increases: a cubic touching the larger-angle ray lies strictly inside the wider Matignon region associated with the smaller angle. Hence
\[
h_{\alpha_1}(b)>h_{\alpha_2}(b)
\qquad
\forall b>0.
\]
Evaluating the two objectives at a minimizer of the \(\alpha_1\) problem gives the stated strict inequality.
\end{proof}

\begin{figure}[t]
\centering
\includegraphics[width=.68\linewidth]{fig_threshold_ratio.pdf}
\caption{Fractional enlargement of Cain's threshold for representative invariant triples. Theorem~\ref{thm:alpha-monotonicity} gives strict decrease in \(\alpha\), and Theorem~\ref{thm:C14} gives convergence to \(T_1\) at integer order. The ordinate is logarithmic.}
\label{fig:threshold-ratio}
\end{figure}

Set
\[
S=\sqrt{\beta_{12}}+\sqrt{\beta_{13}}+\sqrt{\beta_{23}},
\qquad
\Gamma=\sqrt{\beta_{12}\beta_{13}\beta_{23}}.
\]

\begin{theorem}[Sharp classical-limit rate]
\label{thm:C14}
As \(\alpha\to1^-\),
\[
T_\alpha(\beta)
=
T_1(\beta)
+
C(\beta)(1-\alpha)
+
O((1-\alpha)^2),
\]
where
\[
T_1(\beta)=S^2,
\qquad
C(\beta)
=
\pi
\left(
\frac{S^{5/2}}{\sqrt\Gamma}
+
S^{3/2}\sqrt\Gamma
\right).
\]
\end{theorem}

\begin{proof}
Put \(\varepsilon=1-\alpha\). Since
\[
u=\cos(\alpha\pi/2)
=
\frac{\pi}{2}\varepsilon+O(\varepsilon^3),
\]
implicit differentiation of the normalized boundary equation at \(u=0\) gives
\[
h_\alpha(b)
=
b+\pi\varepsilon\sqrt b(1+b)+O(\varepsilon^2).
\]
At \(\alpha=1\) the unique simplex minimizer is
\[
x_1^*=\frac{\sqrt{\beta_{23}}}{S},
\quad
x_2^*=\frac{\sqrt{\beta_{13}}}{S},
\quad
x_3^*=\frac{\sqrt{\beta_{12}}}{S},
\]
with
\[
x_1^*x_2^*x_3^*=\frac{\Gamma}{S^3},
\qquad
B_\beta(x^*)=\frac{\Gamma}{S}.
\]
The envelope theorem, justified by Theorem~\ref{thm:strict-convexity}, therefore gives
\[
C(\beta)
=
\pi
\frac{\sqrt{\Gamma/S}\,(1+\Gamma/S)}
{\Gamma/S^3},
\]
which simplifies to the stated formula.
\end{proof}

\subsection{Minimum dimension for robust genuinely fractional separation}

\begin{theorem}[Minimum robust dimension]
\label{thm:C09}
For every \(0<\alpha<1\),
\[
\min\left\{
n:
\operatorname{int}
\left(
\mathcal F_\alpha^{(n)}
\setminus
\mathcal D_H^{(n)}
\right)
\neq\varnothing
\right\}
=
3.
\]
\end{theorem}

\begin{proof}
For \(n=1\), the two stability notions coincide. For \(n=2\), Corollary~\ref{cor:2x2-empty-interior} gives empty interior.

For \(n=3\), use \(A_\gamma\). It satisfies \(-A_\gamma\) strict P for every \(\gamma>0\), and it is Hurwitz D-stable exactly for \(\gamma<2\). If \(0<\alpha\le2/3\), Kellogg's theorem gives \(A_\gamma\in\mathcal F_\alpha^{(3)}\) for every \(\gamma\), so any \(\gamma>2\) lies in the genuinely fractional class. If \(2/3<\alpha<1\), Corollary~\ref{cor:cyclic-family} gives fractional D-stability for \(\gamma<R_3(\alpha)\), and \(R_3(\alpha)>2\) because
\[
T_\alpha(1,1,1)>T_1(1,1,1)=9.
\]
Choose \(2<\gamma<R_3(\alpha)\). In both regimes the defining inequalities are strict, so the difference class contains an open neighborhood of the chosen matrix.
\end{proof}

The proof-complete technical version of the manuscript also develops a global parametrization of the threshold surface, exact invariant-space sensitivities, the \(\alpha\downarrow2/3\) asymptotic, a closed-form sufficient certificate, and a general realizability theorem. Those secondary results are retained in the accompanying technical manuscript and reproducibility archive but are not needed for the central elimination theorem or the minimum-dimension conclusion.
